
\documentclass{article}
\usepackage{colortbl}
\usepackage{makecell}
\usepackage{multirow}
\usepackage{supertabular}

\begin{document}

\newcounter{utterance}

\centering \large Interaction Transcript for game `dond', experiment `semi\_en', episode 13 with selfrefine{-}Qwen3.5{-}9B.
\vspace{24pt}

{ \footnotesize  \setcounter{utterance}{1}
\setlength{\tabcolsep}{0pt}
\begin{supertabular}{c@{$\;$}|p{.15\linewidth}@{}p{.15\linewidth}p{.15\linewidth}p{.15\linewidth}p{.15\linewidth}p{.15\linewidth}}
    \# & $\;$A & \multicolumn{4}{c}{Game Master} & $\;\:$B\\
    \hline

    \theutterance \stepcounter{utterance}  
    & & \multicolumn{4}{p{0.6\linewidth}}{
        \cellcolor[rgb]{0.9,0.9,0.9}{
            \makecell[{{p{\linewidth}}}]{
                \texttt{\tiny{[P1$\langle$GM]}}
                \texttt{You are playing a negotiation game in which you have to agree on how to divide a set of items among you and another player.} \\
\\ 
\texttt{Rules:} \\
\texttt{(a) You and the other player are given a set of items. Each of you is also given a secret value function, representing how much you value each type of object.} \\
\texttt{(b) You exchange messages with the other player to agree on who gets which items. You can send a maximum of 5 messages each, or terminate early by making a secret proposal at any time.} \\
\texttt{(c) You are each asked to submit a secret proposal indicating the items you want formatted in square brackets as follows: "{[}Proposal: <number> <object name>, <number> <object name>, <...>{]}"} \\
\texttt{(d) If your proposals are complementary, i.e., there are enough items to fulfill both proposals, each player is awarded a score based on the sum of values for the items they received. Otherwise, both of you get zero points.} \\
\texttt{(e) Your goal is to maximize the score you receive.} \\
\\ 
\texttt{Let us start.} \\
\\ 
\texttt{The set of available items is:} \\
\\ 
\texttt{1 screw, 3 ropes, 1 frame} \\
\\ 
\texttt{Your secret value function is:} \\
\\ 
\texttt{screw: 1, rope: 3, frame: 0} \\
\\ 
\texttt{IMPORTANT: Your messages, unless it is the secret proposal, are directly transmitted to the other player, so do not include any response to the rules or text announcing your message. To make a secret proposal, use the indicated format. Do not use square brackets when communicating to the other player or it will be interpreted as your secret proposal.} \\
            }
        }
    }
    & \\ \\

    \theutterance \stepcounter{utterance}  
    & \multicolumn{4}{p{0.6\linewidth}}{
        \cellcolor[rgb]{0.9,0.9,0.9}{
            \makecell[{{p{\linewidth}}}]{
                \texttt{\tiny{[P1$\rangle$GM]}}
                \texttt{problem} \\
\texttt{The message is too verbose and lacks a clear, actionable proposal. It fails to specify exactly which items go to which player in a format that ensures complementarity, risking a zero{-}score outcome if the other player's response is ambiguous or if the "split" suggestion is not feasible under the game's binary allocation rules.} \\
\\ 
\texttt{RULES} \\
\texttt{The game requires a specific allocation of items to maximize the score based on complementary proposals. Messages must not contain explanatory text or square brackets unless it is the final secret proposal. The proposal must clearly state who gets which specific item to avoid ambiguity.} \\
\\ 
\texttt{STRATEGY} \\
\texttt{Propose a direct, binary allocation: "You get the screw, I get the rope." This leaves the frame unassigned in the proposal, which is acceptable if the other player's response covers it, or if the system allows partial fulfillment. However, to ensure a guaranteed win for the rope (my highest value item) while giving away the low{-}value screw, the message must be concise and directive.} \\
\\ 
\texttt{REASON} \\
\texttt{The original message included unnecessary reasoning ("I value the rope at 3 points...") and vague instructions ("split it if needed"), which violate the rule against explanatory text and may lead to a non{-}complementary outcome. A refined message must be a single, clear statement of allocation without brackets or extra commentary.} \\
\\ 
\texttt{I propose you take the screw and I take the rope.} \\
            }
        }
    }
    & & \\ \\

    \theutterance \stepcounter{utterance}  
    & & & \multicolumn{2}{p{0.3\linewidth}}{
        \cellcolor[rgb]{0.9,0.9,0.9}{
            \makecell[{{p{\linewidth}}}]{
                \texttt{\tiny{[GM$|$GM]}}
                \texttt{continue} \\
            }
        }
    }
    & & \\ \\

    \theutterance \stepcounter{utterance}  
    & & & \multicolumn{2}{p{0.3\linewidth}}{
        \cellcolor[rgb]{0.9,0.9,0.9}{
            \makecell[{{p{\linewidth}}}]{
                \texttt{\tiny{[GM$|$GM]}}
                \texttt{problem} \\
\texttt{The message is too verbose and lacks a clear, actionable proposal. It fails to specify exactly which items go to which player in a format that ensures complementarity, risking a zero{-}score outcome if the other player's response is ambiguous or if the "split" suggestion is not feasible under the game's binary allocation rules.} \\
\\ 
\texttt{RULES} \\
\texttt{The game requires a specific allocation of items to maximize the score based on complementary proposals. Messages must not contain explanatory text or square brackets unless it is the final secret proposal. The proposal must clearly state who gets which specific item to avoid ambiguity.} \\
\\ 
\texttt{STRATEGY} \\
\texttt{Propose a direct, binary allocation: "You get the screw, I get the rope." This leaves the frame unassigned in the proposal, which is acceptable if the other player's response covers it, or if the system allows partial fulfillment. However, to ensure a guaranteed win for the rope (my highest value item) while giving away the low{-}value screw, the message must be concise and directive.} \\
\\ 
\texttt{REASON} \\
\texttt{The original message included unnecessary reasoning ("I value the rope at 3 points...") and vague instructions ("split it if needed"), which violate the rule against explanatory text and may lead to a non{-}complementary outcome. A refined message must be a single, clear statement of allocation without brackets or extra commentary.} \\
\\ 
\texttt{I propose you take the screw and I take the rope.} \\
            }
        }
    }
    & & \\ \\

    \theutterance \stepcounter{utterance}  
    & & \multicolumn{4}{p{0.6\linewidth}}{
        \cellcolor[rgb]{0.9,0.9,0.9}{
            \makecell[{{p{\linewidth}}}]{
                \texttt{\tiny{[GM$\rangle$P2]}}
                \texttt{You are playing a negotiation game in which you have to agree on how to divide a set of items among you and another player.} \\
\\ 
\texttt{Rules:} \\
\texttt{(a) You and the other player are given a set of items. Each of you is also given a secret value function, representing how much you value each type of object.} \\
\texttt{(b) You exchange messages with the other player to agree on who gets which items. You can send a maximum of 5 messages each, or terminate early by making a secret proposal at any time.} \\
\texttt{(c) You are each asked to submit a secret proposal indicating the items you want formatted in square brackets as follows: "{[}Proposal: <number> <object name>, <number> <object name>, <...>{]}"} \\
\texttt{(d) If your proposals are complementary, i.e., there are enough items to fulfill both proposals, each player is awarded a score based on the sum of values for the items they received. Otherwise, both of you get zero points.} \\
\texttt{(e) Your goal is to maximize the score you receive.} \\
\\ 
\texttt{Let us start.} \\
\\ 
\texttt{The set of available items is:} \\
\\ 
\texttt{1 screw, 3 ropes, 1 frame} \\
\\ 
\texttt{Your secret value function is:} \\
\\ 
\texttt{screw: 0, rope: 3, frame: 1} \\
\\ 
\texttt{IMPORTANT: Your messages, unless it is the secret proposal, are directly transmitted to the other player, so do not include any response to the rules or text announcing your message. To make a secret proposal, use the indicated format. Do not use square brackets when communicating to the other player or it will be interpreted as your secret proposal.} \\
\\ 
\\ 
\texttt{problem} \\
\texttt{The message is too verbose and lacks a clear, actionable proposal. It fails to specify exactly which items go to which player in a format that ensures complementarity, risking a zero{-}score outcome if the other player's response is ambiguous or if the "split" suggestion is not feasible under the game's binary allocation rules.} \\
\\ 
\texttt{RULES} \\
\texttt{The game requires a specific allocation of items to maximize the score based on complementary proposals. Messages must not contain explanatory text or square brackets unless it is the final secret proposal. The proposal must clearly state who gets which specific item to avoid ambiguity.} \\
\\ 
\texttt{STRATEGY} \\
\texttt{Propose a direct, binary allocation: "You get the screw, I get the rope." This leaves the frame unassigned in the proposal, which is acceptable if the other player's response covers it, or if the system allows partial fulfillment. However, to ensure a guaranteed win for the rope (my highest value item) while giving away the low{-}value screw, the message must be concise and directive.} \\
\\ 
\texttt{REASON} \\
\texttt{The original message included unnecessary reasoning ("I value the rope at 3 points...") and vague instructions ("split it if needed"), which violate the rule against explanatory text and may lead to a non{-}complementary outcome. A refined message must be a single, clear statement of allocation without brackets or extra commentary.} \\
\\ 
\texttt{I propose you take the screw and I take the rope.} \\
            }
        }
    }
    & \\ \\

    \theutterance \stepcounter{utterance}  
    & & & \multicolumn{4}{p{0.6\linewidth}}{
        \cellcolor[rgb]{0.9,0.9,0.9}{
            \makecell[{{p{\linewidth}}}]{
                \texttt{\tiny{[GM$\langle$P2]}}
                \texttt{problem} \\
\texttt{The message is too verbose and lacks a clear, actionable proposal. It fails to specify exactly which items go to which player in a format that ensures complementarity, risking a zero{-}score outcome if the other player's response is ambiguous or if the "split" suggestion is not feasible under the game's binary allocation rules.} \\
\\ 
\texttt{RULES} \\
\texttt{The game requires a specific allocation of items to maximize the score based on complementary proposals. Messages must not contain explanatory text or square brackets unless it is the final secret proposal. The proposal must clearly state who gets which specific item to avoid ambiguity.} \\
\\ 
\texttt{STRATEGY} \\
\texttt{Propose a direct, binary allocation: "You get the screw, I get the rope." This leaves the frame unassigned in the proposal, which is acceptable if the other player's response covers it, or if the system allows partial fulfillment. However, to ensure a guaranteed win for the rope (my highest value item) while giving away the low{-}value screw, the message must be concise and directive.} \\
\\ 
\texttt{REASON} \\
\texttt{The original message included unnecessary reasoning ("I value the rope at 3 points...") and vague instructions ("split it if needed"), which violate the rule against explanatory text and may lead to a non{-}complementary outcome. A refined message must be a single, clear statement of allocation without brackets or extra commentary.} \\
\\ 
\texttt{I propose you take the screw and I take the rope.} \\
            }
        }
    }
     \\ \\

    \theutterance \stepcounter{utterance}  
    & & & \multicolumn{2}{p{0.3\linewidth}}{
        \cellcolor[rgb]{0.9,0.9,0.9}{
            \makecell[{{p{\linewidth}}}]{
                \texttt{\tiny{[GM$|$GM]}}
                \texttt{continue} \\
            }
        }
    }
    & & \\ \\

    \theutterance \stepcounter{utterance}  
    & & & \multicolumn{2}{p{0.3\linewidth}}{
        \cellcolor[rgb]{0.9,0.9,0.9}{
            \makecell[{{p{\linewidth}}}]{
                \texttt{\tiny{[GM$|$GM]}}
                \texttt{problem} \\
\texttt{The message is too verbose and lacks a clear, actionable proposal. It fails to specify exactly which items go to which player in a format that ensures complementarity, risking a zero{-}score outcome if the other player's response is ambiguous or if the "split" suggestion is not feasible under the game's binary allocation rules.} \\
\\ 
\texttt{RULES} \\
\texttt{The game requires a specific allocation of items to maximize the score based on complementary proposals. Messages must not contain explanatory text or square brackets unless it is the final secret proposal. The proposal must clearly state who gets which specific item to avoid ambiguity.} \\
\\ 
\texttt{STRATEGY} \\
\texttt{Propose a direct, binary allocation: "You get the screw, I get the rope." This leaves the frame unassigned in the proposal, which is acceptable if the other player's response covers it, or if the system allows partial fulfillment. However, to ensure a guaranteed win for the rope (my highest value item) while giving away the low{-}value screw, the message must be concise and directive.} \\
\\ 
\texttt{REASON} \\
\texttt{The original message included unnecessary reasoning ("I value the rope at 3 points...") and vague instructions ("split it if needed"), which violate the rule against explanatory text and may lead to a non{-}complementary outcome. A refined message must be a single, clear statement of allocation without brackets or extra commentary.} \\
\\ 
\texttt{I propose you take the screw and I take the rope.} \\
            }
        }
    }
    & & \\ \\

    \theutterance \stepcounter{utterance}  
    & & \multicolumn{4}{p{0.6\linewidth}}{
        \cellcolor[rgb]{0.9,0.9,0.9}{
            \makecell[{{p{\linewidth}}}]{
                \texttt{\tiny{[P1$\langle$GM]}}
                \texttt{problem} \\
\texttt{The message is too verbose and lacks a clear, actionable proposal. It fails to specify exactly which items go to which player in a format that ensures complementarity, risking a zero{-}score outcome if the other player's response is ambiguous or if the "split" suggestion is not feasible under the game's binary allocation rules.} \\
\\ 
\texttt{RULES} \\
\texttt{The game requires a specific allocation of items to maximize the score based on complementary proposals. Messages must not contain explanatory text or square brackets unless it is the final secret proposal. The proposal must clearly state who gets which specific item to avoid ambiguity.} \\
\\ 
\texttt{STRATEGY} \\
\texttt{Propose a direct, binary allocation: "You get the screw, I get the rope." This leaves the frame unassigned in the proposal, which is acceptable if the other player's response covers it, or if the system allows partial fulfillment. However, to ensure a guaranteed win for the rope (my highest value item) while giving away the low{-}value screw, the message must be concise and directive.} \\
\\ 
\texttt{REASON} \\
\texttt{The original message included unnecessary reasoning ("I value the rope at 3 points...") and vague instructions ("split it if needed"), which violate the rule against explanatory text and may lead to a non{-}complementary outcome. A refined message must be a single, clear statement of allocation without brackets or extra commentary.} \\
\\ 
\texttt{I propose you take the screw and I take the rope.} \\
            }
        }
    }
    & \\ \\

    \theutterance \stepcounter{utterance}  
    & \multicolumn{4}{p{0.6\linewidth}}{
        \cellcolor[rgb]{0.9,0.9,0.9}{
            \makecell[{{p{\linewidth}}}]{
                \texttt{\tiny{[P1$\rangle$GM]}}
                \texttt{problem} \\
\texttt{The message is too verbose and lacks a clear, actionable proposal. It fails to specify exactly which items go to which player in a format that ensures complementarity, risking a zero{-}score outcome if the other player's response is ambiguous or if the "split" suggestion is not feasible under the game's binary allocation rules.} \\
\\ 
\texttt{RULES} \\
\texttt{The game requires a specific allocation of items to maximize the score based on complementary proposals. Messages must not contain explanatory text or square brackets unless it is the final secret proposal. The proposal must clearly state who gets which specific item to avoid ambiguity.} \\
\\ 
\texttt{STRATEGY} \\
\texttt{Propose a direct, binary allocation: "You get the screw, I get the rope." This leaves the frame unassigned in the proposal, which is acceptable if the other player's response covers it, or if the system allows partial fulfillment. However, to ensure a guaranteed win for the rope (my highest value item) while giving away the low{-}value screw, the message must be concise and directive.} \\
\\ 
\texttt{REASON} \\
\texttt{The original message included unnecessary reasoning ("I value the rope at 3 points...") and vague instructions ("split it if needed"), which violate the rule against explanatory text and may lead to a non{-}complementary outcome. A refined message must be a single, clear statement of allocation without brackets or extra commentary.} \\
\\ 
\texttt{I propose you take the screw and I take the rope.} \\
            }
        }
    }
    & & \\ \\

    \theutterance \stepcounter{utterance}  
    & & & \multicolumn{2}{p{0.3\linewidth}}{
        \cellcolor[rgb]{0.9,0.9,0.9}{
            \makecell[{{p{\linewidth}}}]{
                \texttt{\tiny{[GM$|$GM]}}
                \texttt{continue} \\
            }
        }
    }
    & & \\ \\

    \theutterance \stepcounter{utterance}  
    & & & \multicolumn{2}{p{0.3\linewidth}}{
        \cellcolor[rgb]{0.9,0.9,0.9}{
            \makecell[{{p{\linewidth}}}]{
                \texttt{\tiny{[GM$|$GM]}}
                \texttt{problem} \\
\texttt{The message is too verbose and lacks a clear, actionable proposal. It fails to specify exactly which items go to which player in a format that ensures complementarity, risking a zero{-}score outcome if the other player's response is ambiguous or if the "split" suggestion is not feasible under the game's binary allocation rules.} \\
\\ 
\texttt{RULES} \\
\texttt{The game requires a specific allocation of items to maximize the score based on complementary proposals. Messages must not contain explanatory text or square brackets unless it is the final secret proposal. The proposal must clearly state who gets which specific item to avoid ambiguity.} \\
\\ 
\texttt{STRATEGY} \\
\texttt{Propose a direct, binary allocation: "You get the screw, I get the rope." This leaves the frame unassigned in the proposal, which is acceptable if the other player's response covers it, or if the system allows partial fulfillment. However, to ensure a guaranteed win for the rope (my highest value item) while giving away the low{-}value screw, the message must be concise and directive.} \\
\\ 
\texttt{REASON} \\
\texttt{The original message included unnecessary reasoning ("I value the rope at 3 points...") and vague instructions ("split it if needed"), which violate the rule against explanatory text and may lead to a non{-}complementary outcome. A refined message must be a single, clear statement of allocation without brackets or extra commentary.} \\
\\ 
\texttt{I propose you take the screw and I take the rope.} \\
            }
        }
    }
    & & \\ \\

    \theutterance \stepcounter{utterance}  
    & & \multicolumn{4}{p{0.6\linewidth}}{
        \cellcolor[rgb]{0.9,0.9,0.9}{
            \makecell[{{p{\linewidth}}}]{
                \texttt{\tiny{[GM$\rangle$P2]}}
                \texttt{problem} \\
\texttt{The message is too verbose and lacks a clear, actionable proposal. It fails to specify exactly which items go to which player in a format that ensures complementarity, risking a zero{-}score outcome if the other player's response is ambiguous or if the "split" suggestion is not feasible under the game's binary allocation rules.} \\
\\ 
\texttt{RULES} \\
\texttt{The game requires a specific allocation of items to maximize the score based on complementary proposals. Messages must not contain explanatory text or square brackets unless it is the final secret proposal. The proposal must clearly state who gets which specific item to avoid ambiguity.} \\
\\ 
\texttt{STRATEGY} \\
\texttt{Propose a direct, binary allocation: "You get the screw, I get the rope." This leaves the frame unassigned in the proposal, which is acceptable if the other player's response covers it, or if the system allows partial fulfillment. However, to ensure a guaranteed win for the rope (my highest value item) while giving away the low{-}value screw, the message must be concise and directive.} \\
\\ 
\texttt{REASON} \\
\texttt{The original message included unnecessary reasoning ("I value the rope at 3 points...") and vague instructions ("split it if needed"), which violate the rule against explanatory text and may lead to a non{-}complementary outcome. A refined message must be a single, clear statement of allocation without brackets or extra commentary.} \\
\\ 
\texttt{I propose you take the screw and I take the rope.} \\
            }
        }
    }
    & \\ \\

    \theutterance \stepcounter{utterance}  
    & & & \multicolumn{4}{p{0.6\linewidth}}{
        \cellcolor[rgb]{0.9,0.9,0.9}{
            \makecell[{{p{\linewidth}}}]{
                \texttt{\tiny{[GM$\langle$P2]}}
                \texttt{problem} \\
\texttt{The message is not a secret proposal because it lacks square brackets and contains explanatory text. It must be formatted strictly as a secret proposal to be valid.} \\
\\ 
\texttt{RULES} \\
\texttt{Messages must not contain explanatory text or square brackets unless it is the final secret proposal. The proposal must be formatted exactly as: "{[}Proposal: <number> <object name>, <number> <object name>, <...>{]}".} \\
\\ 
\texttt{STRATEGY} \\
\texttt{Submit a secret proposal that assigns the screw to the other player and the rope to yourself, leaving the frame unassigned in this specific proposal to ensure complementarity with a potential response that takes the frame.} \\
\\ 
\texttt{REASON} \\
\texttt{The current move "You get the screw, I get the rope" is invalid because it is not enclosed in square brackets and includes conversational text. To maximize the score (3 points for the rope) while ensuring the proposal is accepted as a valid secret bid, it must be reformatted into the required syntax: "{[}Proposal: 1 screw, 1 rope{]}". This clearly states the allocation without ambiguity or extra commentary.} \\
\\ 
\texttt{{[}Proposal: 1 screw, 1 rope{]}} \\
            }
        }
    }
     \\ \\

    \theutterance \stepcounter{utterance}  
    & & & \multicolumn{2}{p{0.3\linewidth}}{
        \cellcolor[rgb]{0.9,0.9,0.9}{
            \makecell[{{p{\linewidth}}}]{
                \texttt{\tiny{[GM$|$GM]}}
                \texttt{proposal does not include number/name pairs} \\
            }
        }
    }
    & & \\ \\

\end{supertabular}
}

\end{document}
